\centerline{\bf \huge 10 класс}
\centerline{\bf \huge Решения}

\problem{10.1}
Корни уравнения $x^2+p x+q=0$ - целые числа, отличные от нуля. Докажите, что число $p^2+(q-1)^2$ составное.

\textit{Решение.}
Пусть $a$ и $b$ - корни данного уравнения. Тогда, используя теорему Виета, получим: $p^2+(q-1)^2=(a+b)^2+(a b-1)^2=a^2+2 a b+b^2+a^2 b^2-2 a b+1=\left(a^2+1\right)+\left(b^2+\right. \left.a^2 b^2\right)=\left(a^2+1\right)\left(b^2+1\right)$. Так как $a \neq 0$ и $b \neq 0$, то $a^2+1 \neq 1$ и $b^2+1 \neq 1$. Следовательно, число $p^2+(q-1)^2$ составное.

\problem{10.2}
Можно ли квадрат со стороной 8 полностью покрыть двумя кругами диаметра 9?

\textit{Ответ: можно.}

\textit{Решение.}
Разделим квадрат на два равных прямоугольника размером $4 \times 8$. Диагональ такого прямоугольника равна $\sqrt{4^2+8^2}=4 \sqrt{5}<9$. Следовательно, каждый из полученных прямоугольников можно покрыть описанным около него кругом, диаметр которого равен длине диагонали прямоугольника. Значит, и кругом с диаметром 9 такой прямоугольник покрыть можно. Тогда данный квадрат можно покрыть двумя кругами диаметра 9.

\problem{10.3}
Пусть $n$ - натуральное число. Какая цифра стоит сразу после запятой в десятичной записи числа $\sqrt{n^2+n}$ ?

\textit{Ответ: 4.}

\textit{Решение.}
Докажем, что $(n+0,4)^2<n^2+n<(n+0,5)^2$, откуда и будет следовать указанный ответ. Действительно:

1) $(n+0,4)^2<n^2+n \Leftrightarrow 0,8 n+0,16<n \Leftrightarrow n>0,8$. Последнее неравенство верно для любого натурального $n$.

2) $n^2+n<(n+0,5)^2 \Leftrightarrow n<n+0,25$, а это очевидно.

\problem{10.4}
Вписанная окружность треугольника $A B C$ касается стороны $A C$ в точке $D$. Вторая окружность проходит через точку $D$, касается луча $B A$ в точке $A$ и, кроме того, касается продолжения стороны $B C$ за точку $C$. Найдите отношение $A D: D C$.

\textit{Ответ: $3: 1$.}

\textit{Решение.} Пусть вторая окружность касается луча $B C$ в точке $K$, а вписанная окружность касается сторон $A B$ и $B C$ треугольника в точках $P$ и $Q$ соответственно (см. рис. 10.4). Далее можно рассуждать по-разному.

\textbf{Первый способ}. 
Пусть $C D=C Q=x$ и $A D=A P= kx$. Так как $B A=B K$ и $B P=B Q$, то $P A=Q K$, откуда $C K =k x-x$. По теореме о касательной и секущей: $CD\cdot CA =C K^2$, то есть $x(x+k x)=(k x-x)^2$. Разделив обе части этого равенства на $x^2$, раскрывая скобки и приводя подобные слагаемые, получим: $k^2=3 k$, то есть $k=3$.

\textbf{Второй способ.}
Пусть $B C=a, A C=b, A B=c$, тогда $C D=\dfrac{a+b+c}{2}-c=\dfrac{a+b-c}{2}, C K= B K-B C=A B-B C=c-a$. По теореме о касательной и секущей: $C D \cdot C A=C K^2$, то есть $\dfrac{a+b-c}{2} \cdot b=(c-a)^2 \Leftrightarrow b^2-(c-a) b-2(c-a)^2=0$. Решая это уравнение как квадратное относительно $b$, получим: $b=a-c$ или $b=2(c-a)$.

Равенство $b=a-c$ противоречит неравенству треугольника, а если $b=2(c-a)$, то $\dfrac{A D}{D C}=\dfrac{b+c-a}{2}: \dfrac{a+b-c}{2}=\dfrac{3 c-3 a}{c-a}=3$.


\problem{10.5}
У царя восемь сыновей, и все дураки. Каждую ночь царь отправляет троих из них стеречь золотые яблоки от жар-птицы. Поймать жар-птицу царевичи не могут, винят в этом друг друга, и поэтому никакие двое не соглашаются пойти вместе в караул второй раз. Какое наибольшее количество ночей это может продолжаться?


\textit{Ответ:8 ночей.}

\textit{Решение.}
\textit{Оценка.} Рассмотрим любого из сыновей. В каждом карауле он находится вместе с двумя братьями, поэтому после трёх его выходов останется один брат, вместе с которым он не ходил в караул, и уже не сможет пойти, так для них не будет третьего. Эта ситуация «симметрична», то есть если сын А после трёх своих выходов не ходил в караул с сыном $B$, то и сын В после трёх своих выходов не ходил в караул с сыном $A$. Следовательно, есть как минимум четыре пары сыновей, которые в одном карауле не были и не будут. Всевозможных пар: $\frac{8 \cdot 7}{2}=28$, поэтому в караулах побывают не более чем $28-4=24$ пары. Так как за одну ночь задействованы по три пары сыновей, то ночей может быть не больше чем $24: 3=8$.

\textit{Пример.} Обозначим сыновей: $1, 2, 3, 4, A, B, C, D.$ Тогда восемь троек, удовлетворяющих условию задачи: $A12, B23, C34, D41, AB4, BC1, CD2, DA3.$

%Муницип Москва 2019 10 класс